\subsection{Semiprimordial Modules}

DEFINITION 4. --- A module is called semiprimordial if it is the direct sum of a family of primordial submodules.

*Examples. --- 1) Every simple module is primordial (VIII, p.~45); every semisimple module is therefore semiprimordial (VIII, p.~55, Definition 1).

\begin{enumerate}
\def\labelenumi{\arabic{enumi})}
\setcounter{enumi}{1}
\tightlist
\item
  If A is a left Noetherian ring, then every injective A-module is semiprimordial (X, §1, n° 9, p.~21, proposition 14 and X, §1, n° 10, p.~22, théorème 3, b)).*
\end{enumerate}

THEOREM 1 (Azumaya). --- Let A be a ring, L a primordial A-module, and M a semiprimordial A-module. There exists a unique cardinal, denoted by {[}M : L{]}, with the following property:

For every decomposition \(M = \bigoplus_{i \in I} M_i\) of M as a direct sum of primordial modules, the set of indices \(i \in I\) such that \(M_i\) is isomorphic to L has cardinal {[}M : L{]}.

The proof is based on the following four lemmas.

Lemma 1. --- Let M be an A-module, M' a primordial submodule of M, and M'\,' a submodule of M supplementary to M'. Let u be an endomorphism of M. Then u or \(1_{M} - u\) induces an isomorphism from M' to a submodule of M supplementary to M'\,'.

Let \(p\) be the projection from M onto \(M'\) with kernel \(M''\) , and let \(v\) be the restriction of \(p \circ u\) to \(M'\) . First, suppose that \(v\) is an automorphism of \(M'\) . Since \(v\) is injective, the restriction of \(u\) to \(M'\) is injective, and we have \(u(M') \cap M'' = 0\) . Since \(v\) is surjective, we have \(u(M') \oplus M'' = M\) . Consequently, \(u\) induces an isomorphism from \(M'\) to a submodule supplementary to \(M''\) in M. Now, suppose that \(v\) is not an automorphism of \(M'\) . Then \(1_{M'} - v\) is an automorphism of \(M'\) because \(M'\) is primordial. Now, \(1_{M'} - v\) is the restriction of \(p \circ (1_M - u)\) to \(M'\) . The previous reasoning proves that \(1_M - u\) induces an isomorphism from \(M'\) to a submodule of M supplementary to \(M''\) .

Lemma 2. --- Let M be an A-module that is the direct sum of a family \((\mathrm{M}_{i})_{i\in\mathrm{I}}\) of primordial submodules, and let u be an endomorphism of M. Set \(v=1_{M}-u\) and \(M_{J}=\bigoplus_{i\in J}M_{i}\) for every subset J of I. Then one of the following two properties holds:

\begin{enumerate}
\def\labelenumi{\alph{enumi})}
\item
  There exists an index \(i \in I\) such that u induces an isomorphism from \(M_{i}\) to a direct factor submodule of M.
\item
  For every finite subset \(J\) of \(I\) , \(v\) induces an isomorphism from \(M_J\) to a submodule supplementary to \(M_{I-J}\) .
\end{enumerate}

If property b) holds, then v is injective.

Suppose that property a) does not hold, and let us establish property b) by induction on the cardinal of J. There is nothing to prove if \(J = \varnothing\) . Therefore, assume that J is nonempty, choose an element i of J, and set \(J' = J - \{i\}\) . By the induction hypothesis, v induces an isomorphism from \(M_{J'}\) to a submodule of M supplementary to \(M_{I-J'} = M_{I-J} \oplus M_i\) ; consequently, the submodule \(M'' = v(M_{J'}) \oplus M_{I-J}\) is supplementary to \(M_i\) . By Lemma 1 and the assumption on u, the endomorphism v induces an isomorphism from \(M_i\) to a submodule of M supplementary to \(M''\) ; consequently, v induces an isomorphism from \(M_J = M_i \oplus M_{J'}\) to a submodule supplementary to \(M_{I-J}\) .

The last assertion follows from the fact that M is the union of the submodules \(M_{J}\) , where J runs through the finite subsets of I.

Lemma 3. --- Let M be an A-module that is the direct sum of a family \((\mathrm{M}_{i})_{i\in\mathrm{I}}\) of primordial submodules, and let p be a nonzero projector of M. There exists an index \(i\in I\) such that p induces an isomorphism from \(M_{i}\) to a direct factor submodule of \(p(\mathrm{M})\) .

Since p is nonzero, \(1_{M} - p\) is not injective. By Lemma 2, there exists an index \(i \in I\) such that p induces an isomorphism from \(M_{i}\) to a direct factor submodule of M. Every projector of M with image \(p(\mathrm{M}_{i})\) defines, by restriction, a projector of \(p(\mathrm{M})\) with image \(p(\mathrm{M}_{i})\) , so \(p(\mathrm{M}_{i})\) is a direct factor submodule of \(p(\mathrm{M})\) .

Lemma 4. --- Let M be an A-module that is the direct sum of a family \((\mathrm{M}_{i})_{i\in\mathrm{I}}\) of primordial submodules, let L be a primordial A-module, and let N be a direct factor submodule of M. Suppose that N is the direct sum of a family \((\mathrm{N}_{j})_{j\in\mathrm{J}}\) of submodules isomorphic to L, and denote by \(I_{L}\) the set of indices \(i\in I\) such that \(M_{i}\) is isomorphic to L. We then have

\[
\operatorname{Card} (\mathrm{J}) \leqslant \operatorname{Card} \left(\mathrm{I} _ {\mathrm{L}}\right).\tag{1}
\]

Let \(N_0\) be a submodule of M supplementary to N. The module M is the direct sum of \(N_0\) and the family \((N_j)_{j \in J}\) . For every \(j \in J\) , denote by \(p_j\) the projector of M with image \(N_j\) associated with this decomposition (II, §1, No.~8, p.~209, Proposition 12). For every \(i \in I\) , denote by \(J(i)\) the set of indices \(j \in J\) such that \(p_j\) induces an isomorphism from \(M_i\) to \(N_j\) . This set is finite: indeed, if \(x\) is a nonzero element of \(M_i\) and K is a finite subset of J such that \(x\) belongs to \(\mathrm{N}_0 + \sum_{k\in \mathrm{K}}\mathrm{N}_k\) , then we have \(p_j(x) = 0\) for \(j\in \mathrm{J} - \mathrm{K}\) , so that \(\mathrm{J}(i)\) is contained in \(\mathrm{K}\) .

Let \(j \in J\) . By Lemma 3, there exists an index \(i \in I\) such that \(p_j\) induces an isomorphism from \(M_i\) to a direct factor submodule of \(N_j\) . Since \(M_i\) is nonzero and \(N_j\) is primordial and therefore indecomposable (VIII, p.~31, Proposition 4), we have \(p_j(M_i) = N_j\) , and \(j\) belongs to \(J(i)\) . Since the module \(M_i\) is isomorphic to \(N_j\) and hence to \(L\) , the index \(i\) belongs to \(I_L\) . This proves that \(J\) is the union of the family of finite sets \((J(i))_{i \in I_L}\) . If the set \(J\) is infinite, then the set \(I_L\) is infinite, and we have (Set Theory, III, §6, No.~3, p.~188, Corollary 3)

\[
\operatorname{Card} (\mathrm{J}) \leqslant \sum_ {i \in \mathrm{I} _ {\mathrm{L}}} \operatorname{Card} (\mathrm{J} (i)) \leqslant \operatorname{Card} (\mathrm{I} _ {\mathrm{L}}).
\]

Now, suppose that the set J is finite, and let us prove the lemma by induction on the cardinal of J. If J is empty, there is nothing to prove. Therefore, suppose that J is nonempty, and choose an element j of J. By the above, there exists an index \(i \in I_{L}\) such that \(p_{j}\) induces an isomorphism from \(M_{i}\) to \(N_{j}\) . Set \(I' = I - \{i\}\) and \(J' = J - \{j\}\) . The module M is the direct sum of \(M_{i}\) and the kernel of \(p_{j}\) . It is also the direct sum of \(M_{i}\) and the submodule \(M' = \oplus_{i' \in I'} M_{i'}\) . Hence, there exists (II, §1, No.~9, p.~210, Corollary of Proposition 13) an isomorphism \(\varphi\) from \(\operatorname{Ker} p_{j} = N_{0} \oplus_{j' \in J'} N_{j'}\) to \(M'\) . Set \(N' = \varphi\big(\sum_{j' \in J'} N_{j'}\big)\) . The submodule \(N'\) of \(M'\) is a direct factor and is the direct sum of the family \((\varphi(N_{j'}))_{j' \in J'}\) of primordial submodules isomorphic to L. Let us apply the induction hypothesis to \(M'\) and \(N'\) : we have \(\text{Card}(J') \leqslant \text{Card}(I_L - \{i\})\) and therefore inequality (1).

Let us prove Theorem 1. Let \((\mathrm{M}_{i})_{i\in\mathrm{I}}\) and \((\mathrm{N}_{j})_{j\in\mathrm{J}}\) be two families of primordial submodules with direct sum M. Let \(I_{L}\) (resp. \(J_{L}\) ) be the set of \(i\in I\) (resp. \(j\in J\) ) such that \(M_{i}\) (resp. \(N_{j}\) ) is isomorphic to L. We have \(\mathrm{Card}(J_{\mathrm{L}})\leqslant\mathrm{Card}(I_{\mathrm{L}})\) by Lemma 4. By interchanging the roles of I and J, we obtain the inverse inequality and thus the theorem.

The cardinal \([M : L]\) defined in Theorem 1 is called the primordial multiplicity of L in M.

COROLLARY 1. --- Let M and N be semiprimordial modules. Then M and N are isomorphic if and only if we have \([M:L]=[N:L]\) for every primordial module L.

COROLLARY 2. --- Let M be a semiprimordial module. Let \((\mathrm{M}_{i})_{i\in\mathrm{I}}\) and \((\mathrm{M}_{j}^{\prime})_{j\in\mathrm{J}}\) be families of primordial submodules of M such that

\[
\mathrm{M} = \bigoplus_ {i \in \mathrm{I}} \mathrm{M} _ {i} = \bigoplus_ {j \in \mathrm{J}} \mathrm{M} _ {j} ^ {\prime}.
\]

Then there exist an automorphism \(u\) of \(M\) and a bijection \(\varphi\) from \(I\) to \(J\) such that we have \(u(M_i) = M_{\varphi(i)}'\) for every \(i \in I\) .

For any primordial module L, let \(I_{L}\) (resp. \(J_{L}\) ) be the set of indices \(i \in I\) (resp. \(j \in J\) ) such that \(M_{i}\) (resp. \(M_{j}'\) ) is isomorphic to L. The nonempty sets of the form \(I_{L}\) (resp. \(J_{L}\) ) form a partition of I (resp. J), and for every L, we have

\[
\mathrm{Card} (\mathrm {I_ {L}}) = \mathrm{Card} (\mathrm {J_ {L}}) = [ \mathrm{M:L} ];
\]

the corollary follows.

COROLLARY 3. --- Let M, N, and P be semiprimordial modules. Suppose that \(M \oplus P\) is isomorphic to \(N \oplus P\) and that \([P : L]\) is finite for every primordial module L. Then M and N are isomorphic.

By assumption, we have

\[
[ \mathrm{M}: \mathrm{L} ] + [ \mathrm{P}: \mathrm{L} ] = [ \mathrm{N}: \mathrm{L} ] + [ \mathrm{P}: \mathrm{L} ]
\]

for every primordial module L. Since \([P:L]\) is finite, it follows by induction from (Set Theory, III, §3, No.~4, p.~162, Proposition 8) that we have \([M:L]=[N:L]\) for every primordial module L. The modules M and N are therefore isomorphic by Corollary 1.

COROLLARY 4. --- Let M and N be semiprimordial modules. Suppose that there exists an integer d \textgreater{} 0 such that \(M^{d}\) is isomorphic to \(N^{d}\) . Then the modules M and N are isomorphic.

Let \(\mathrm{L}\) be a primordial module. By assumption, we have

\[
d [ \mathrm{M}: \mathrm{L} ] = d [ \mathrm{N}: \mathrm{L} ].
\]

We therefore have the equality \([M:L]=[N:L]\) : indeed, we have da = a for every infinite cardinal a (Set Theory, III, §6, No.~3, p.~188, Corollary 4). The modules M and N are then isomorphic by Corollary 1.

COROLLARY 5. --- Let M be a semiprimordial module that is the direct sum of a finite family \((\mathrm{M}_{i})_{i\in\mathrm{I}}\) of primordial submodules. For any subset J of I, write \(M_{J}=\bigoplus_{i\in J}M_{i}\) . Let N be a direct factor submodule of M.

\begin{enumerate}
\def\labelenumi{\alph{enumi})}
\item
  There exists a subset \(\mathrm{J}\) of \(\mathrm{I}\) such that \(\mathrm{M}_{\mathrm{J}}\) is a submodule supplementary to \(\mathrm{N}\) .
\item
  Let \(\mathrm{J}\) be a subset of \(\mathrm{I}\) . If \(\mathrm{M}_{\mathrm{J}}\) is supplementary to \(\mathrm{N}\) , then the module \(\mathrm{N}\) is isomorphic to \(\mathrm{M}_{\mathrm{I - J}}\) and is semiprimordial.
\end{enumerate}

We denote by K the set of indices \(i \in I\) such that \(N \cap M_{i} = 0\) ; let us use induction on the cardinal of K. The corollary is clear if M = N. Suppose \(M \neq N\) . Let p be a projector of M with kernel N. Denote its image by P. It is nonzero, and by Lemma 3, there exists a \(j \in I\) such that p induces an isomorphism from \(M_{j}\) to a direct factor submodule of P. We have \(N \cap M_{j} = 0\) . Set \(N' = N \oplus M_{j}\) . We have \(N' = N \oplus p(M_{j})\) . A submodule supplementary to \(p(M_{j})\) in P is also supplementary to \(N'\) in M, so that \(N'\) is a direct factor submodule of M. The set of indices \(i \in I\) such that \(N' \cap M_{i} = 0\) is contained in \(K - \{j\}\) . By the induction hypothesis, there exists a subset \(J'\) of I such that \(M_{J'}\) is a submodule supplementary to \(N'\) in M. Set \(J = J' \cup \{j\}\) . Then \(M_{J}\) is a submodule supplementary to N in M.

Let J be a subset of I such that \(M_{J}\) is a submodule supplementary to M in N. Since \(M_{J}\) is also supplementary to \(M_{I-J}\) , the modules N and \(M_{I-J}\) are isomorphic and N is semiprimordial.

COROLLARY 6. --- Every finitely generated projective module over a local ring is free. \(^{(2)}\)

Let A be a local ring. The A-module \(A_{s}\) is primordial (VIII, p.~31). If M is a finitely generated projective A-module, then there exist an A-module N and a natural number n such that \(M \oplus N\) is isomorphic to \(A_{s}^{n}\) (II, §2, No.~2, p.~232, Corollary 1). It follows from Corollary 5 that the module M is itself isomorphic to \(A_{s}^{m}\) for an integer m such that \(0 \leqslant m \leqslant n\) , hence is free.

Remark. --- Let M and \(M'\) be semiprimordial A-modules. It immediately follows from Lemma 4 of VIII, p.~33 that \(M'\) is isomorphic to a direct factor submodule of M if and only if we have \([M':L] \leqslant [M:L]\) for every primordial A-module L. In particular, if L is a primordial A-module, \([M:L]\) is the greatest of the cardinals a for which there exists a direct factor submodule of M isomorphic to \(\mathrm{L}^{(\mathfrak{a})}\) .

\pandocbounded{\includegraphics[keepaspectratio]{images/part001_3d6df7c1.jpg}}

The relation \([M:L]=0\) therefore means that there exists no direct factor submodule of M isomorphic to L. This does not exclude the existence of a submodule of M isomorphic to L; it suffices to consider the example where A=Z, L=Z/2Z, and M=Z/4Z: the Z-modules L and M are primordial and not isomorphic, so that \([M:L]=0\) and L is isomorphic to the submodule \(2Z/4Z\) of M.

